\chapter{Design and Implementation}
\markboth{Design and Implementation}{Design and Implementation}

\section{UWB}
In swarm drone projects, one information is central and will determine the achievement of the formation shaping.
This information is the \textbf{position}. In our case, we are measuring the relative position between drone. \\ 
To achieve this, our professor imposed UWB as localisation sensors. Those sensors are based on electromagnetic signal and needed two modules to interact with each others (Find more information in the bibliographic study \ref{Bibliographical Study}.

\subsection{Algorithm choice}
As exposed in the bibliographic study, there are different algorithms which respond to multiple needs. In this project, we need precise localisation data and no interference, furthermore we have a maximun 3 drones in the swarm. \\
With those caracterestics in mind, we choose the DS-TWR location algorithm (\ref{DS-TWR}) and the TDMA timing algorithm (\ref{TDMA}). \\
We have implemented them into the ESP32-DW1000 module. The $\href{https://github.com/M1chigan/Projet-S-IoT-Essaim-de-Drone/blob/main/Code/Code_commun/src/uwb_driver.cpp}{code}$ is using C/C++ as the software language. 

\subsection{Calibration}

In general, sensors need a calibration to use them efficiently. For UWB module, we need to establish a distance reference. For this project, the reference is the expected distance between drones: 1.20m. Afterwards, we had to adapt the antenna delay parameters which represent the signal's time spent in the antenna's PCB. It's important to note that each module has a different antenna time delay. To calibrate, we had to choose a module as a reference, here called UWB 1, and then modify the antenna delay of each drone to make them a 1.20m measurement match exactly with 1.20m in real life. After this step, all the drones' module are calibrated to work with each other. After doing some tests, we decided to modify the reference to 1.30m, this modification is explained in the next part \ref{FtFV}.
\newpage
\subsection{UWB Test}
 During last year's project, our colleagues found out that UWB module should be installed vertically to optimize their precision and reach. Therefore, we have repeated the same test in order to verify their result. 
\subsubsection{Face to Face verticaly}
\label{FtFV}
We have carried out same vertical test following this scheme :
\begin{figure}[H]
  \centering
  \FtFVertical
  \caption{Head-to-head vertical configuration}
\end{figure}

\begin{figure}[H]
  \centering
  \begin{minipage}[c]{0.48\textwidth}
    \centering
    \footnotesize
    \setlength{\tabcolsep}{3pt}
    \resizebox{\linewidth}{!}{%
    \begin{tabular}{|l|l|l|}
    \hline
        Real distance (cm) & Measured distance (cm) & Error  \\ \hline
        70 & 67 & 3 \\ \hline
        80 & 78 & 2 \\ \hline
        90 & 89,2 & 0,8 \\ \hline
        100 & 100 & 0 \\ \hline
        110 & 110,6 & -0,6 \\ \hline
        120 & 120 & 0 \\ \hline
        130 & 132,4 & -2,4 \\ \hline
        140 & 143,6 & -3,6 \\ \hline
        150 & 154,2 & -4,2 \\ \hline
        160 & 165,7 & -5,7 \\ \hline
        170 & 176,4 & -6,4 \\ \hline
    \end{tabular}}
    \caption{Measures for head-to-head vertical modules, referenced at 120cm}
  \end{minipage}\hfill
  \begin{minipage}[c]{0.48\textwidth}
    \centering
    \footnotesize
    \setlength{\tabcolsep}{3pt}
    \resizebox{\linewidth}{!}{%
    \begin{tabular}{|l|l|l|}
    \hline
        Real distance (cm) & Measured distance (cm) & Error  \\ \hline
        70 & 64 & 6 \\ \hline
        80 & 76,2 & 3,8 \\ \hline
        90 & 85,6 & 4,4 \\ \hline
        100 & 97,4 & 2,6 \\ \hline
        110 & 108,2 & 1,8 \\ \hline
        120 & 119,6 & 0,4 \\ \hline
        130 & 130 & 0 \\ \hline
        140 & 140 & 0 \\ \hline
        150 & 152,7 & -2,7 \\ \hline
        160 & 162 & -2 \\ \hline
        170 & 172,2 & -2,2 \\ \hline
    \end{tabular}}
    \caption{Measures for head-to-head vertical modules, referenced at 130cm}
  \end{minipage}
\end{figure}
\noindent As we can observe in these tables, the accuracy is decreasing fast when the module is getting further from the reference so placing the reference at 130cm would allow us to get better accuracy at $\pm$ 10cm around 120cm.
\subsubsection{Face to Face horizontaly}
\begin{figure}[H]
  \centering

  % --- Schéma ---
  \resizebox{0.5\linewidth}{!}{\FtFhorizontale}
  \captionof{figure}{Head-to-head horizontal configuration}

  \vspace{1em}

  % --- Tableau en deux moitiés ---
  \small
  \setlength{\tabcolsep}{4pt}
  \begin{minipage}[t]{0.46\textwidth}
    \vspace{0pt}
    \centering
    \begin{tabular}{|c|c|c|}
      \hline
      Real distance (cm) & Measured distance (cm) & Error \\ \hline
      69,4  & 61    & -8,4 \\ \hline
      70    & 62,2  & -7,8 \\ \hline
      75    & 69,6  & -5,4 \\ \hline
      80    & 73,9  & -6,1 \\ \hline
      82,6  & 76,1  & -6,5 \\ \hline
      85    & 79,9  & -5,1 \\ \hline
      90    & 86,6  & -3,4 \\ \hline
      91    & 87    & -4   \\ \hline
      95    & 92,2  & -2,8 \\ \hline
      100   & 96,5  & -3,5 \\ \hline
      100   & 96,8  & -3,2 \\ \hline
      105   & 102,6 & -2,4 \\ \hline
      107,3 & 106,4 & -0,9 \\ \hline
      110   & 108,9 & -1,1 \\ \hline
      115   & 113,4 & -1,6 \\ \hline
      120,5 & 120,9 & 0,4  \\ \hline
      125   & 126,4 & 1,4  \\ \hline
      130   & 129,6 & -0,4 \\ \hline
      132   & 133,7 & 1,7  \\ \hline
    \end{tabular}
  \end{minipage}\hfill
  \begin{minipage}[t]{0.46\textwidth}
    \vspace{0pt}
    \centering
    \begin{tabular}{|c|c|c|}
      \hline
      Real distance (cm) & Measured distance (cm) & Error \\ \hline
      135   & 133,6 & -1,4 \\ \hline
      140   & 138,6 & -1,4 \\ \hline
      145   & 144,9 & -0,1 \\ \hline
      147,4 & 150   & 2,6  \\ \hline
      150   & 151,7 & 1,7  \\ \hline
      155   & 160,3 & 5,3  \\ \hline
      155,5 & 161   & 5,5  \\ \hline
      160   & 165,7 & 5,7  \\ \hline
      164,7 & 171,8 & 7,1  \\ \hline
      165   & 172,2 & 7,2  \\ \hline
      167   & 172,8 & 5,8  \\ \hline
      170   & 177,2 & 7,2  \\ \hline
      175   & 180,3 & 5,3  \\ \hline
      176   & 183,8 & 7,8  \\ \hline
      180   & 184,3 & 4,3  \\ \hline
      185   & 188,3 & 3,3  \\ \hline
      189,6 & 197,9 & 8,3  \\ \hline
      190   & 191,9 & 1,9  \\ \hline
    \end{tabular}
  \end{minipage}

  \captionof{table}{Measures in the head-to-head horizontal configuration}
\end{figure}

Let's compare both positions (horizontal or vertical) : as we can observe around 120cm for horizontal position, we got about 1.8cm of error and for vertical, we got about 1cm. Also, when the distance moves off 120 cm, the accuracy decreases slower in vertical than in horizontal. Based on this experimentation, we chose to use our module vertically to get better accuracy and flexibility. \\

\noindent Normally, in most radio module documentations, there is an antenna diagram which describes the efficient area of the antenna. We were looking for this diagram in order to choose the position of the drone which gets the best UWB accuracy. We did not find it, therefore we had to do some tests in different configurations.

\subsubsection{Parallel modules}

\begin{figure}[H]
  \centering
  \begin{minipage}[c]{0.38\textwidth}
    \centering
    \resizebox{\linewidth}{!}{\SchemaParalleles}
    \captionof{figure}{Parallel configuration}
  \end{minipage}\hfill
  \begin{minipage}[c]{0.58\textwidth}
    \centering
    \footnotesize
    \setlength{\tabcolsep}{3pt}
    \resizebox{\linewidth}{!}{%
    \begin{tabular}{|c|c|c|c|c|}
      \hline
      offset (cm) & Angle (°) & Real distance (cm) & Measured distance (cm) & Error \\ \hline
      0 & 0 & 120 & 117 & 3 \\ \hline
      10 & 4,78 & 121 & 118,4 & 2,6 \\ \hline
      20 & 9,59 & 121,5 & 118,9 & 2,6 \\ \hline
      30 & 14,48 & 122,5 & 119,1 & 3,4 \\ \hline
      40 & 19,47 & 125,8 & 124,8 & 1 \\ \hline
      50 & 24,62 & 128,5 & 128,7 & -0,2 \\ \hline
      60 & 30,00 & 132,5 & 134,4 & -1,9 \\ \hline
      70 & 35,69 & 137 & 142,6 & -5,6 \\ \hline
    \end{tabular}}
    \captionof{table}{Measures in the parallel configuration}
  \end{minipage}
\end{figure}


\subsubsection{Radius configuration}
\begin{figure}[H]
  \centering
  \begin{minipage}[c]{0.38\textwidth}
    \centering
    \resizebox{\linewidth}{!}{\Schemarayon}
    \captionof{figure}{Radius configuration}
  \end{minipage}\hfill
  \begin{minipage}[c]{0.58\textwidth}
    \centering
    \footnotesize
    \setlength{\tabcolsep}{3pt}
    \resizebox{\linewidth}{!}{%
    \begin{tabular}{|l|l|l|}
    \hline
        Angle(°) & Real distance (cm) & Measured distance (cm)  \\ \hline
        10 & 120 & 116,7 \\ \hline
        20 & 120 & 118,4 \\ \hline
        30 & 120 & 119,5 \\ \hline
        40 & 120 & 122,8 \\ \hline
        50 & 120 & 125,9 \\ \hline
        60 & 120 & 124,2 \\ \hline
        70 & 120 & 123 \\ \hline
        80 & 120 & 124 \\ \hline
        90 & 120 & 130 \\ \hline
    \end{tabular}}
    \captionof{table}{Measures in the radius configuration}
  \end{minipage}
\end{figure}

\subsubsection{Head-to-back configuration}
\begin{figure}[H]
  \centering
  \begin{minipage}[c]{0.38\textwidth}
    \centering
    \resizebox{\linewidth}{!}{\FtB}
    \captionof{figure}{Head-to-back configuration}
  \end{minipage}\hfill
  \begin{minipage}[c]{0.58\textwidth}
    \centering
    \footnotesize
    \setlength{\tabcolsep}{3pt}
    \resizebox{\linewidth}{!}{%
    \begin{tabular}{|l|l|l|}
    \hline
        Real distance (cm) & Measured distance (cm) & Error  \\ \hline
        70 & 97 & -27 \\ \hline
        80 & 109,8 & -29,8 \\ \hline
        90 & 122,8 & -32,8 \\ \hline
        100 & 133,3 & -33,3 \\ \hline
        110 & 143,5 & -33,5 \\ \hline
        120 & 152 & -32 \\ \hline
        130 & 161,6 & -31,6 \\ \hline
        140 & 177,2 & -37,2 \\ \hline
        150 & 188 & -38 \\ \hline
        160 & 192,2 & -32,2 \\ \hline
        170 & 208,8 & -38,8 \\ \hline
    \end{tabular}}
    \captionof{table}{Measures in the Head-to-back configuration}
  \end{minipage}
\end{figure}

As an overall conclusion to our different configuation tests, we can say that the final distance errors are satisfying with an average error lower than 5cm for every configuration besides one. The so-called "head-to-back" configuration is showing very bad results with an average error above 30cm which is not acceptable for us. This shows that it will be necessary for us to make the different modules face each other to get the best possible results. 

\newpage
\subsection{Filter}
Electromagnetics signals are sensitive to exterior noise and can be interfered easily. Therefore, we have to filter the raw data. To do so we can implement three different filters.

\subsubsection{Polynomial filter}
During calibration and tests, we have discovered that the errors are linked to the real distance and vary following an approximate polynomial equation. 

\begin{figure}[H]
   \centering
   \resizebox{0.6\textwidth}{!}{\includegraphics{Polynomial.png}}
   \caption{Diagram of the error based on real distance}
\end{figure}

\noindent Using the Excel linearization tool, we got this error equation: 

\begin{equation}
   f(x) = -1{,}004 \times 10^{-6}\, x^4 + 0{,}0005059\, x^3 - 0{,}09263\, x^2 + 7{,}398\, x - 221{,}7
\end{equation}

Every measure passes through the equation and is recalibrated based on this equation. That allows to reduce the error when we are at more than $\pm$ 10 cm from the reference.\\

\noindent We have implemented it in C/C++ as follows: 
\begin{figure}[H]
   \centering
   \resizebox{0.7\textwidth}{!}{\includegraphics{Polynomial_Filter.png}}
   \caption{\href{https://github.com/M1chigan/Projet-S-IoT-Essaim-de-Drone/blob/main/Code/Code_commun/src/filtre.cpp}{Code} for polynomial filter}
\end{figure}
\newpage
\subsubsection{NLOS Filter}
Electromagnetic signals are very sensitive to obstacles. After doing research, we have discovered that the signal will lose power if it encounters an obstacle. Based on that, we can verify that the power of the received signal is strong enough otherwise the value cannot be considered as viable so we can deny it.\\

\noindent We have implemented it in C/C++ as follows: 

\begin{figure}[H]
   \centering
   \resizebox{0.70\textwidth}{!}{\includegraphics{NLOS Filter.png}}
   \caption{\href{https://github.com/M1chigan/Projet-S-IoT-Essaim-de-Drone/blob/main/Code/Code_commun/src/filtre.cpp}{Code} for NLOS Filter}
\end{figure}


\subsubsection{Innovation Gate Filter}
UWB has a tendency to be noisy which can create absurd values, either very high or very low. This data can affect the corrector and create some discontinuity in the drone control. Based on the maximum speed of the drone, we can determine absurdity by calculating the difference between the previous distance $d_{-1}$ and the current $d$, then determine if the difference can be reached at the maximum speed, if not the measure is ignored, otherwise, we accept it.\\

\noindent We have implemented it in C/C+: 

\begin{figure}[H]
   \centering
   \resizebox{0.70\textwidth}{!}{\includegraphics{Innovation_Gate.png}}
   \caption{\href{https://github.com/M1chigan/Projet-S-IoT-Essaim-de-Drone/blob/main/Code/Code_commun/src/filtre.cpp}{Code} for innovation gate filter}
\end{figure}

\newpage

\section{Communication}
In the project, we have two different use for communication: first, communication is used to make each drone talk with each other to maintain the shape but also adapt themselves to the others. Secondly, communication helps the data go from the drone to the base to permit visualisation of the swarm.

\subsection{Communication Protocol Choice}
We chose ESP-NOW as the communication protocol for two main reasons : it doesn't need new hardware implementation and it responds to both communication purposes of the project. Then, as mentioned during the presentation, the code has been optimized and written by the manufacturer which facilitate its implementation. For exemple, sending data requires only this command :
\begin{figure}[H]
   \centering
   \resizebox{0.4\textwidth}{!}{\includegraphics{ESP-NOW_send.png}}
   \caption{Code to send messsages using ESP-NOW}
\end{figure}

As said before, the range of ESP-NOW reaches 150 meters which allows us to create a base. It only requires us to have a ESP32 and connect it to our computer.

\subsection{ESP-NOW implementation}
To implement ESP-NOW, we had to think about the data we want to exchange. For our first test, we are using two different packages for two different purposes. \\
First for inter-drone communication : 
\begin{figure}[H]
   \centering
   \resizebox{0.4\textwidth}{!}{\includegraphics{Drone_package.png}}
   \caption{Drone communication package}
\end{figure}

And for the base to drone communcation, we have this package : 

\begin{figure}[H]
   \centering
   \resizebox{0.4\textwidth}{!}{\includegraphics{base_package.png}}
   \caption{Drone communication package}
\end{figure}


\newpage

\section{Measure processing and formation control}

After studying the UWB measurements and the communication with ESP-NOW, we have to consider the processing of the retrieved data to finally implement an optimal formation control algorithm. The questions that need to be asked are the following : what data do we need to retrieve and what algorithm do we use for its processing and for the formation control. As mentioned before, our drones include many onboard sensors and benefit from the added UWB module. The bibliographical study gave us insights about potential architectures of measure and control so a choice has to be made considering the hardware available. Here are two things to notice :


\begin{itemize}
    \item UWB modules return purely relative distance measurements and give no sense of direction at all
    \item The onboard Extended Kalman Filter from the flight controller's firmware Ardupilot fuses the measurements from its different sensors (accelerometer, gyroscope, magnetometer and barometer) and the added GNSS module to get notably an accurate estimation of the drone's attitude (yaw, roll, pitch) as a quaternion, 3 NED (North East Down) velocity components and 3 NED position components (up to 22 states are estimated in total).
    
\end{itemize}

 Therefore, we know that we can get good estimates of these states through the onboard EKF of the flight controller. A first choice has to be made about what we want to control. To maximize the use of UWB technology, we would want to implement a distance control architecture so that each drone tries to keep a determined distance between itself and other drones. The following steps for the architecture choice are going to, first, figure out if it possible to keep a formation while moving the swarm as a whole then to study different formation control algorithms and finally to decide the final shape of the swarm (for example, 3 drones pointing at the North or 3 drones pointing towards the center of the formation).


\newpage

\section{Visualisation}
In outdoor tests, drones are flying at more than 10 meters away from us. That constraint reduces our visibility and comprehension of the drone movements, attitude and position. To prevent this issue, we need to develop a visualisation system. That would increase our error detection capability and give us the opportunity to save tests data and simulate them afterward to analyse and detect problematic points. Visualisation would also allow us to observe in real time differents parameters such as: attitude, 3D shape positionning, and also RAW data coming from the drone.In order to achieve this system, we chose the Matlab software as a visualizer and software language.  

\subsection{Matlab}

Matlab is a general tool developped by MathWorks. It has many functions such as mathematics calculator, simulator, and also a corrector designer. Based on add-ons and toolboxes, it can be useful for a large extent of purposes. It also allows to generate an interface. 

\subsection{Visualisation Architecture}

The visualisation follows the following architecture : 
\begin{figure}[H]
  \centering
   \resizebox{0.9\textwidth}{!}{\SchVisualisation}
  \caption{Visualization architecture}
\end{figure}

Let's explain each different block and its role : 

\begin{itemize}
    \item ESP32 : Relay between ESP-NOW protocol and the UART Port of the computer, which allows to make it visible for the Matlab code.
    \item COM PORT: Physical bridge between the computer and the ESP32
    \item mainMenu : General window which redirects to the next window
    \item runObserver : Window to observe UART communication on the PORT (Debug utility)
    \item runSimulation : Main Window where we visualise attitude of the drone, shape placement and the data reader
\end{itemize}

\newpage
\section{Simulation}
\label{Simulation}

Due to logistical problem, outdoor tests are rare, and INSA Strasbourg does not have an indoor test room. Howewer, our work need to be tested beforehand and confronted to the reality. To achieve this, we had to develop a simulation protocol which is the as close as possible to the reality. \\
To do so, we had to choose the same firmware as our flight controler and an UAV simulation tool. \\
In this part, we will first present both the firmware and ground station software : respectively ArduPilot and Mission planner, and secondly, explain the simulation architecture, then present the result we obtained and finally the simulation possibilities.

\subsection{ArduPilot and Mission Planner}

To do drone simulation, we have to select the same architecture as our drone with the Pixhawk 4 flight controller. In the drone industry, two main different firmwares exist: Ardupilot and PX4. In our case, the Pixhawk 4 is using the ArduPilot firmware. \\
Then, after selecting the right firmware, we have to use a software to visualise and simulate drones. We chose Mission Planner because we had already worked with it on a past project and it is dedicated to Ardupilot vehicles. \\
Let's present both tools:

\subsubsection{\href{https://github.com/ArduPilot/ardupilot}{ArduPilot}}

 ArduPilot was created in 2007 by the DIY Drones community. Originally developed for Arduino boards, it has since become a general-purpose, open-source autopilot. It can control many types of vehicles (copters, planes, rovers, submarines, etc.). ArduPilot provides autopilot functions, receives data from sensors, and fuses it using an Extended Kalman Filter (EKF) [ref]. It also allows users to create pre-planned missions. To operate, it uses the MAVLink protocol to communicate with the drone and facilitate control. In our case, an ESP32 uses the MAVLink library to send MAVLink messages that command the drone’s movements.

\subsubsection{\href{https://ardupilot.org/planner/}{Mission planner}}

Mission planner is a representative tool for copters such as aerial drones. It was created by Michael Oborne to help users of the ArduPilot architecture to visualise their drone and command them. The software also allows to declare multiple drones and so, create a swarm. For our project, it will be used as a simulation field. 
\newpage
\subsection{Architecture}

The simulation works with a very specific architecture because tools like ArduPilot work under the Linux Operating System. The architecture is exposed in the figure below. 

\begin{figure}[H]
  \centering
   \resizebox{0.9\textwidth}{!}{\SchSimulation}
  \caption{Simulation architecture}
\end{figure}

Let's explain each different bloc and its role : 
\begin{itemize}
    \item $ESP-32_{Base}$ : Represents order sent from the base (e.g. takeoff order) 
    \item $ESP-32_{Drone}$: Represents the main drone where the code we want to test is running
    \item ArduPilot : Represents the Pixhawk's firmware of our drone. Here, we are simulating 3 of them. 
    \item \href{https://github.com/M1chigan/Projet-S-IoT-Essaim-de-Drone/blob/main/Simulation/Script%20d'essai/Bridge.py}{Bridge.py} : Code made to redirect data from the UART connection to UDP connection 
    \item MAVProxy : Manages all inputs and outputs, allows to link the drone inputs and redirect them to Ardupilot and also allows to redirect outputs from ArduPilot (GPS, speed, etc) to Mission Planner.
    \item Mission Planner : Used to visualise our swarm drone based on data from ArduPilot
    \item \href{https://github.com/M1chigan/Projet-S-IoT-Essaim-de-Drone/blob/main/Simulation/Script%20d'essai/Simulation.py}{Simulation.py}  : Python script to simulate $Drone_2$ and $Drone_3$
\end{itemize}


\newpage
\section{Tests}
Tests are the main parts of a technical project as reality is always different from the theory. Therefore, we have to carry out them through two different environements: simulation and outdoor.
\subsection{Simulation}
Using the architecture described in the previous part, we did and planned to do different tests for different parts of the project.  
\subsubsection{Ongoing : }
\label{Simu-level1}
\underline{Console control:} \\
This test has for purpose to test MAVLink functions, communication and drone control.\\ 
\noindent It follows the following scheme: \\

\begin{figure}[H]
   \centering
   \resizebox{0.7\linewidth}{!}{\LevelUN}
   \caption{Architecture of the first level of simulation}
\end{figure}

\noindent The drone is for now capable of these functions (lexicon below) : 

\begin{figure}[H]
   \centering
   \resizebox{0.3\textwidth}{!}{\includegraphics{lexique.png}}
   \caption{Drone capability at level 1}
\end{figure}

\noindent We will be able to conclude that our communication between base and drone is well-designed and that the MAVLink procedure is successfully sent to the drone. 

\subsubsection{Planned :}

\noindent  \underline{Level 2:} \\ 
\label{Simu-level2}
The objective of this level is : a slave drone needs to follow the master drone controlled by the console as in the level 1. To do so it should use communication, attitude and GPS data. Possibly, we can think of integrating UWB to make a full test, closer to the real situation.   

\begin{figure}[H]
   \centering
   \resizebox{0.8\linewidth}{!}{\LevelDeux}
   \caption{Architecture of the second level of simulation}
\end{figure}

\noindent \underline{Level 3:} \\ 
The objective of this level is to reach 3 drones and introduce a drone swarm. Out of the 3 drones, two of them would be controlled by a Python script. The third drone is the tested one. It allows to test the different correctors. 

\begin{figure}[H]
   \centering
   \resizebox{0.8\linewidth}{!}{\LevelTrois}
   \caption{Architecture of the third level of simulation}
\end{figure}

\newpage
\noindent \underline{Level 4:} \\
Objective of this level is to reach a full simulation environement where all the three drone are working on it own ESP32. Mission planer and Ardupilot are strictly an environement simulation. It is the closest simulation to the objective of the project.

\begin{figure}[H]
   \centering
   \resizebox{0.8\linewidth}{!}{\LevelQuatre}
   \caption{Architecture of the fourth level of simulation}
\end{figure}

\subsection{Outdoor}
Make drone fly outdoors is restricted by the law. It makes test difficult to proceed and rare. Therefore, outdoor tests needs to be precisely prepared to not lose any time and be efficient on site.
\subsubsection{Planned}
We have planned to attempt outdoor tests within 2 weeks. During this test, we would like to test: 
\begin{itemize}
  \item UWB placement (rotation, and height)
  \item ESP-NOW range 
  \item Drone control console (as simulation level 1 [\ref{Simu-level1}])
  \item Slave drone following master drone (as simulation level 1 [\ref{Simu-level2}])
\end{itemize} 


